\scanpage{359}
taking traces we find
\[
  2\lambda=2 \tag{45}
\]

--- these are the necessary and sufficient conditions on $\lambda$ for
\deleted{\unsure{$\lambda$}}
($\varepsilon$ being regular unipotent)
$\lambda\varepsilon$ to be still \inserted{\unsure{strict\textsuperscript{ly}}}
\unsure{unipoten\textsuperscript{t}} (hence necessari\textsuperscript{ly} regular unipotent).
If $2$ is invertible on $S$, \inserted{\uline{more generally}
if $2$ is a regular section of $\underline{\mathcal O}_S$,}
\deleted{\unsure{is reduced}}
these conditions imply $\lambda=1$, OK, but
\deleted{\unsure{in the case}}\inserted{outside} of this hypothesis (for example in characteristic $2$
\deleted{\unsure{even if}} over a nonreduced scheme) this is no longer necessarily so. Let us write
\[
  \lambda=1+\nu
\]
the conditions (44) (45) read
\[
  2\nu=0,\quad \nu^2=0, \tag{46}
\]
thus for the map (43) to be an injection on $S$, it is necessary and sufficient
that the \inserted{only} sections of $S$ which satisfy (46)
are the zero section.

In general, denoting by $\mu'_2$ the subgroup scheme
of $\mu_2$ defined by (45)
(it is in fact an absolute infinitesimal group \inserted{scheme},
\deleted{\unsure{over}} \inserted{finite over $\mathbb Z$, not flat over $\mathbb Z$,}
isomorphic to the spectrum of
\[
  \mathbb Z[T]/(T^2-1,2T-2)\simeq
  \mathbb Z[S]/(S^2,2\unsure{S}),
\]
also isomorphic to the scheme of dual numbers over $\mathbb Z$,
relative to the module $\mathbb Z/2\mathbb Z$ over $\mathbb Z$,
\deleted{\unsure{$\mu'_2\simeq$}}
\[
  \mu'_2\simeq\operatorname{Spec}
  \bigl(\mathbb Z\deleted{\unsure{\varepsilon}}\oplus\underbrace{\mathbb Z/2\mathbb Z}_{\text{square-zero ideal}}\bigr)
  \quad\ldots
\]
\mspage{358}

% Printed archive footer outside the manuscript frame:
% 358
% https://grothendieck.umontpellier.fr
% Université de Montpellier, UMR CNRS 5149, Case courrier 051, Place Eugène Bataillon
% 34095 Montpellier cedex 5 - France

\scanpage{360}
\mspage{481}
It turns out that $\mu'_{2S}$ operates
\inserted{freely by translations}
on the scheme of regular unipotents \inserted{\unsure{(or not)}}
of $G$, and we have
\[
\diagram{\begin{tikzcd}[ampersand replacement=\&]
\underline{\mathrm{Unip.rég}}(G)/\mu'_{2S}
  \arrow[r,"\sim"] \& \underline{\mathrm{Unip.rég}}(PG)
\end{tikzcd}}\ . \tag{47}
\]
This seems at first sight to be in contradiction with the \unsure{sequence} of Borels
\[
\diagram{\begin{tikzcd}[ampersand replacement=\&,column sep=small,row sep=large,nodes in empty cells]
\deleted{\unsure{1}} \& B \arrow[d,mapsto]
 \& \underline{\mathrm{Borels}}(G)
   \arrow[r,"\sim","B\mapsto B/\mathbb \mathbb G_m"']
   \arrow[d,"\wr"']
 \& \mathrm{Borels}(PG) \arrow[d,"\wr"]
 \& \deleted{\unsure{B}}\widetilde B \arrow[d,mapsto] \\
 \& B_u
 \& \begin{gathered}\text{``Subgroups}\\
     \deleted{\unsure{\mathrm U}}\text{unipotent with 1}\\
     \text{parameter'' of }G\end{gathered}
     \arrow[r,"\sim"]
 \& \begin{gathered}\text{``Subgroups}\\
     \text{unipotent with 1 parameter''}\\\text{of }PG\end{gathered}
 \& \widetilde B_u \\
 \& \& \text{\unsure{that is}}\quad L \arrow[r,mapsto,"\sim"] \& p(L)
\end{tikzcd}} \tag{48}
\]
(where $p:G\to PG$ is the canonic\textsuperscript{al} morphism)

Could the difficulty come from the fact that the vertical arrows
(surjective, by definition of the notion of ``1-parameter unipotent subgroups''
of $G$ or of $PG$) might not be injective, outside the hypothesis that $2$ is invertible
(i.e. \unsure{residue} characteristic\ $\ne2$)?
I feel like saying that
\[
  \underline B=\operatorname{Norm}(B_u),\quad
  \widetilde B=\operatorname{Norm}(\widetilde B_u) \tag{49}
\]
are formulas which allow us to recover $B,\widetilde B$ with the help
of $B_u,\widetilde B_u$\unsure{,|''} --- and I have
\deleted{\unsure{done}} indeed \uline{\unsure{already}}
checked the first of these formulas, by dint of linear algebra in $M$
(in the form: the 1-parameter unipotent subgroups of $G$
correspond $1$-$1$ to the sections of $\mathbb P$ over $S$,
i.e. to the line \unsure{subbundles} $D$ of $F$\ldots) ---
\mspage{359}

% Printed archive footer outside the manuscript frame:
% 359
% https://grothendieck.umontpellier.fr
% Université de Montpellier, UMR CNRS 5149, Case courrier 051, Place Eugène Bataillon
% 34095 Montpellier cedex 5 - France

\scanpage{361}
but \unsure{this same} point in $PG$, it is
\deleted{if $u\in\Gamma M^*$ is such that
$u\widetilde N u^{-1}=\mu N+\lambda$,}
\inserted{not very trivial.}
\deleted{with $\mu\in\Gamma\mathbb \mathbb G_m$,}

The difficulty must rather come from the fact \unsure{that a|that no}
\unsure{regular unipotent} section $\widetilde\varepsilon$ of $PG$
does not define (without restriction on $S$) a 1-parameter unipotent subgroup
--- whereas for a section $\varepsilon$ of
\unsure{regular unipotents|regular unipotent} of $G$
\[
  \varepsilon=1+N,
\]
we define \inserted{one} by the formula
\[
  \varepsilon(t)=1+tN\ .
\]
We \unsure{will note} that if we replace $\varepsilon$ by
$\inserted{\varepsilon'=}\lambda\varepsilon$
(with $\lambda$ satisfying (44) (45) $\lambda^2=1$, $2\lambda=2$),
We have
\[
  \varepsilon'=(1+\nu)(1+N)=1+N'
  \qquad N'=(\unsure{1+\nu})N+\nu
\]
and
\[
  \varepsilon'(t)=1+tN'
\]
I say that if \inserted{$\lambda\ne1$ i.e.}
$\nu\ne0$, \deleted{\unsure{we then have}} the composite homomorphisms
\[
\diagram{\begin{tikzcd}[ampersand replacement=\&]
\widetilde{\underline\varepsilon},\widetilde{\underline\varepsilon}'\colon
\mathbb G_{\unsure{a|0}} \arrow[r] \& G \arrow[r] \& PG
\end{tikzcd}}
\]
\margin{$\Vert$}
\deleted{\unsure{will be}} associated to $\varepsilon,\varepsilon'$
will be \uline{distinct}, although they have
the same value for $t=1$.
Indeed, \deleted{\unsure{we will have}} the derivatives at the origin
of $\widetilde{\underline\varepsilon},\widetilde{\underline\varepsilon}'$
are identified with the composites
\[
\diagram{\begin{tikzcd}[ampersand replacement=\&]
\underline{\mathcal O}_S \arrow[r] \& M \arrow[r] \& M/\underline{\mathcal O}_S\unsure{=|\simeq}P\unsure{\mathfrak g}
\end{tikzcd}}
\]
\mspage{360}

% Printed archive footer outside the manuscript frame:
% 360
% https://grothendieck.umontpellier.fr
% Université de Montpellier, UMR CNRS 5149, Case courrier 051, Place Eugène Bataillon
% 34095 Montpellier cedex 5 - France

\scanpage{362}
\mspage{482}
which send $1$ to $\widetilde N$ resp. to $\widetilde N'$,
(where the $\sim$ here denotes passage to the \uline{additive} quotient) we have
\[
  \widetilde N'=(1+\nu)\widetilde N
\]
and $\widetilde N$ is a section of $P\unsure{\mathfrak g}$ everywhere nonzero
--- thus we can have \deleted{not}
\[
  \widetilde N'=\widetilde N\quad\text{i.e.}\quad
  (1+\nu)\widetilde N=\widetilde N\quad\text{i.e.}\quad
  \nu\widetilde N=0
\]
only if $\nu=0$, i.e. $\lambda=1+\nu=1$, q.e.d.

In conclusion, we must consider that we give ourselves in $PG$,
not only $\widetilde\rho$ and $\widetilde\sigma$
(whence $\widetilde\varepsilon_0=\widetilde\sigma^{-1}\widetilde\rho$,
$\widetilde\varepsilon_1=\widetilde\rho\widetilde\sigma^{-1}$),
but also ``1-parameter subgroups'' (parametrized)
\[
\diagram{\begin{tikzcd}[ampersand replacement=\&,nodes in empty cells]
\deleted{\mathbb G_{\unsure{a|0}}}
  \arrow[r,"\deleted{\unsure{\widetilde{\underline\varepsilon}_0}}"] \& \deleted{PG}
\end{tikzcd}}\ ,
\]
\[
\left\{\diagram{\begin{tikzcd}[ampersand replacement=\&,row sep=small]
\widetilde{\underline\varepsilon}_0\colon \mathbb G_{\unsure{a|0}}
  \arrow[r] \& PG \\
\widetilde{\underline\varepsilon}_1\colon \mathbb G_{\unsure{a|0}}
  \arrow[r] \& PG
\end{tikzcd}}\right.\ . \tag{49}
\]
related of course by
\[
\left\{\begin{aligned}
\widetilde{\underline\varepsilon}_1
 &=\rho(\widetilde{\underline\varepsilon}_0)
   \bigl(=\sigma(\widetilde{\underline\varepsilon}_0)\bigr)
   \quad\text{i.e.}\\
\widetilde{\underline\varepsilon}_1(t)
 &=\rho\bigl(\widetilde{\underline\varepsilon}_0(t)\bigr)
   =\sigma\bigl(\widetilde{\underline\varepsilon}_0(\unsure{t|u})\bigr)
\end{aligned}\right. \tag{50}
\]
which means precisely that we \uline{give} ourselves, in addition to
$\widetilde\rho,\widetilde\sigma$, \unsure{a} lifting\deleted{\illegible{1}}
\inserted{\inserted{regular} \unsure{unipoten\textsuperscript{t}}}
$\varepsilon_0$ of $\widetilde\varepsilon_0$
(\deleted{\unsure{which|that}} the choice is therefore not \unsure{a} priori unique),
and we deduce $\varepsilon_1$ from it by
\[
  \deleted{\unsure{\rho}}\quad
  \varepsilon_1=\widetilde\rho(\varepsilon_0)
  =\widetilde\sigma(\varepsilon_0)
\]
\margin{we will necessarily have a relation
$\sigma^{-1}\rho= -\lambda_0\varepsilon_0$
hence \deleted{$\unsure{\operatorname{int}(\varepsilon_0)|\mathrm{i.e.}\,\sigma(\varepsilon_0)}\,\unsure{\varepsilon_0}$}
$\unsure{\sigma(\varepsilon_0)|\sigma^{-1}(\varepsilon_0)}=\rho(\varepsilon_0)\,\unsure{=}$}
\mspage{361}

% Printed archive footer outside the manuscript frame:
% 361
% https://grothendieck.umontpellier.fr
% Université de Montpellier, UMR CNRS 5149, Case courrier 051, Place Eugène Bataillon
% 34095 Montpellier cedex 5 - France

\scanpage{363}
\mspage{4}
\deleted{\unsure{And}}
Thus, we have adjusted the \deleted{\unsure{precision|position}}
\inserted{description} of the data, so as to
\deleted{\unsure{adjust}} be able to give precise normalization
statements --- namely

\uline{Proposition} Let $\widetilde\rho,\widetilde\sigma$
be sections of $PG$, and $\widetilde{\underline\varepsilon}_0$
a ``strict'' \inserted{additive} subgroup \unsure{with} 1 parameter
\margin{(i.e. a monomorphism and whose image is a $\widetilde B_u$)}
\[
\diagram{\begin{tikzcd}[ampersand replacement=\&]
\widetilde{\underline\varepsilon}_0\colon
\mathbb G_{\unsure{a|0}}
  \arrow[r,phantom,"\unsure{\longrightarrow|\text{---}}" description] \& PG
\end{tikzcd}}\ .
\]
\unsure{such|such} that we have \unsure{$\times$|$*$}
\[
  \widetilde\sigma^{-1}\widetilde\rho
   =\widetilde{\underline\varepsilon}_0(1)
   \overset{\mathrm{dfn}}=\widetilde\varepsilon_0 \tag{51}
\]
whence
\[
\left\{\begin{aligned}
\widetilde\sigma(\widetilde{\underline\varepsilon}_0)
 &=\widetilde\rho(\widetilde{\underline\varepsilon}_0)
   \overset{\mathrm{dfn}}=\widetilde{\underline\varepsilon}_1,
   \quad\text{and}\\
\widetilde\rho\widetilde\sigma^{-1}
 &=\widetilde{\underline\varepsilon}_1(1)
   \overset{\mathrm{dfn}}=\widetilde\varepsilon_1
   =\rho(\widetilde\varepsilon_0)=\sigma(\widetilde\varepsilon_0)
\end{aligned}\right. \tag{52}
\]

Then:

1) (\unsure{``For the record''})
$\exists$ unique $\varepsilon_0,\varepsilon_1$
\inserted{\deleted{\unsure{strict\textsuperscript{ly}}} regular unipotent} sections
\deleted{such} of $G$ such that
(setting $\varepsilon_i=\deleted{\unsure{\mathrm{Id}+}}1+N_i$,
$\underline\varepsilon_i(t)=1+tN_i$) we have
\[
  \widetilde{\underline\varepsilon}_i(t)=\varepsilon_i(t)^{\sim}\ . \tag{53}
\]

2) \deleted{Suppose that $\varepsilon_0$}
\inserted{We will then have}
\[
  \unsure{\dot{\varepsilon}_1|\varepsilon_1}=\widetilde\sigma(\varepsilon_0)
    =\widetilde\rho(\varepsilon_0) \tag{54}
\]
\[
\deleted{
  \unsure{\rho}\sigma^{-1}\rho=-\lambda_0\varepsilon_0
  \quad\varepsilon_0=-\lambda_0\sigma^{-1}\rho,\quad
  \varepsilon_1=-\lambda_0\rho\sigma^{-1},
}\LMcancelledtag{\deleted{55}}
\]
\deleted{3) For $\varepsilon_0$ and $\varepsilon_1$ to be
everywhere independent \unsure{in} each fiber, i.e. $N$}

3) \unsure{Let|Being} $\rho,\sigma$ be liftings of
$\widetilde\rho,\widetilde\sigma$, we will then have
\[
  \deleted{\unsure{\rho\sigma^{-1}\rho}}\qquad
  \varepsilon_0=-\lambda_0\sigma^{-1}\rho,\quad
  \varepsilon_1=-\lambda_0\rho\sigma^{-1} \tag{55}
\]
for $\lambda_0\in\Gamma\mathbb \mathbb G_m$ uniquely determined.
For a suitable choice of $\rho$ (when $\sigma$ is already chosen)
we can assume $\lambda_0=1$
(which determines $\rho$ uniquely in
\deleted{\unsure{function}} terms of $\sigma$, \unsure{and} conversely)
\mspage{362}

% Printed archive footer outside the manuscript frame:
% 362
% https://grothendieck.umontpellier.fr
% Université de Montpellier, UMR CNRS 5149, Case courrier 051, Place Eugène Bataillon
% 34095 Montpellier cedex 5 - France

\scanpage{364}
\mspage{483}

\unsure{4)} so that we have
\[
\tag{56}
\varepsilon_0=-\unsure{\sigma}^{-1}\rho
\qquad(\text{hence }\varepsilon_1=-\rho\sigma^{-1})
\]

4) $\rho$ and $\sigma$ being chosen so as to satisfy (56), for $\varepsilon_0,\varepsilon_1$ to be \deleted{everywhere} independent \inserted{(in each fiber)}, i.e. for us to have $N_0$ and $N_1$ nonhomothetic in \unsure{every} fiber, it is necessary and sufficient that we have
\[
\tag{57}
\operatorname{Tr}(\rho+\sigma)\quad\text{invertible}
\]
(cf. also proposition page 478' for other equivalent conditions).

5) For $\varepsilon_0,\varepsilon_1$ to be independent \inserted{(in each fiber)}, it is necessary and sufficient that we \deleted{have} can find \inserted{(locally \uline{\unsure{fppf}})} $\rho,\sigma$ above $\widetilde\rho,\widetilde\sigma$ so as to satisfy (56) \uline{and}
\[
\tag{57'}
\operatorname{Tr}(\rho+\sigma)=1
\qquad\text{i.e.}\quad
\operatorname{Tr}\rho+\operatorname{Tr}\sigma=1
\]
\margin{$\rho+\sigma$ an elementary projector $\Longleftrightarrow$}
and this determines $\rho,\sigma$ \uline{uniquely} (in terms of $\widetilde\rho,\widetilde\sigma,\underline{\widetilde\varepsilon}_0$).

\uline{Remarks} 1) The conditions on $(\rho,\sigma)$ considered in the proposition on page 478', \unsure{and} the normalization $\unsure{u|v_1}=1$ i.e. $\operatorname{Tr}(\sigma+\rho)=1$ i.e. $\sigma+\rho$ elementary projector, are \uline{symmetric} in $(\rho,\sigma)$: they are satisfied for $(\rho,\sigma)$ iff they are for $(\sigma,\rho)$ --- \deleted{\unsure{the interchange of $\rho$ and $\sigma$|the interchange is made}} \inserted{to be more} precise let us put
\[
\tag{58}
\rho'=\sigma,\ \sigma'=\rho\qquad
\varepsilon'_0=\underbrace{-\sigma'^{-1}\rho'}_{\unsure{=}-\rho^{-1}\sigma},\quad
\varepsilon'_1=\underbrace{-\rho'\sigma'^{-1}}_{-\sigma\rho^{-1}}
\]
then we will have
\[
\tag{59}
\varepsilon'_0=\varepsilon_0^{-1},\qquad
\varepsilon'_1=\varepsilon_1^{-1}.
\]
Likewise, $(\rho^{-1},\sigma^{-1})$ satisfies the same conditions \inserted{on page 478'}, and \unsure{putting|taking}

\mspage{363}

% Printed archive footer outside the manuscript frame:
% 363
% https://grothendieck.umontpellier.fr
% Université de Montpellier, UMR CNRS 5149, Case courrier 051, Place Eugène Bataillon
% 34095 Montpellier cedex 5 - France

\scanpage{365}

\[
\tag{60}
\rho''=\rho^{-1},\ \sigma''=\sigma^{-1},\qquad
\varepsilon''_0=\underbrace{-\sigma''^{-1}\rho''}_{-\sigma\rho^{-1}},\quad
\varepsilon''_1=\underbrace{-\rho''\sigma''^{-1}}_{-\rho^{-1}\sigma}
\]
we will have
\[
\tag{61}
\varepsilon''_0=\varepsilon_{\unsure{1|0}}^{-1},\qquad
\varepsilon''_1=\varepsilon_0^{-1},
\]
but we will take note that in general we will no longer have $\operatorname{Tr}(\rho''+\unsure{\sigma''})=1$, but (by an immediate matrix calculation)
\[
\tag{62}
\operatorname{Tr}(\rho''+\sigma'')
=\operatorname{Tr}(\rho^{-1}+\sigma^{-1})
=-\frac{1}{\unsure{v|\nu}}
=(\det\rho)^{-1}=(\det\sigma)^{-1}
\]
(assuming of course \inserted{already} $\unsure{v_1|u_1}=1$ i.e. $\operatorname{Tr}(\rho+\sigma)=1$)
--- this is equal to 1 only if
\[
\tag{63}
\det\rho\ (=\det\sigma)=1
\]
--- this is the condition of the \unsure{section|relation} 3$^\circ$) of $\unsure{(A1)_2|(41)_2}$ (page 463'), which therefore also means that the normalization $\unsure{v_1|u_1}=1$ is stable under the transformation $\rho\mapsto\rho^{-1},\ \sigma\mapsto\sigma^{-1}$.

On the other hand, the choice of the \uline{\unsure{pinning}} associated with $(\rho,\sigma)$ (\unsure{i.e.} $\widetilde\rho,\widetilde\sigma$) [in which the emphasis has been put on $\varepsilon_1=-\rho\sigma^{-1}$, which receives the simple form
\[
\varepsilon_1=\begin{pmatrix}1&0\\1&1\end{pmatrix},
\]
whereas $\varepsilon_0$ \inserted{$=-\sigma^{-1}\rho$} has the more complicated form
\[
\varepsilon_0=\begin{pmatrix}1&\unsure{v_1/v|u_1/v}\\0&1\end{pmatrix}
\qquad\left(=\begin{pmatrix}1&-1\\0&1\end{pmatrix}\text{ in the modular case}\right)
\]
] implicitly brings in the order in which we give ourselves $\rho,\sigma$ --- if we change it, we change the \unsure{pinning} by the change from $\varepsilon_{\unsure{1}}$ to $\varepsilon'_1=\varepsilon_1^{-1}$.

In fact, the situation \unsure{offered} the choice between \uline{four} ``\unsure{natural}'' \unsure{pinnings}, associated with the pair of

\mspage{364}

% Printed archive footer outside the manuscript frame:
% 364
% https://grothendieck.umontpellier.fr
% Université de Montpellier, UMR CNRS 5149, Case courrier 051, Place Eugène Bataillon
% 34095 Montpellier cedex 5 - France

\scanpage{366}
\mspage{484}

one-parameter subgroups $\{L_0,L_1\}$, \deleted{and} \inserted{\unsure{corresponding}} to the four sections $\varepsilon_0,\varepsilon_0^{-1},\varepsilon_1,\varepsilon_1^{-1}$ of the one \deleted{\unsure{or}} \inserted{and of} the other.

We preferred to normalize with the help of $L_1$, because of the most convenient form of the loop equation (which I finally arrived at moreover after some rather heavy \unsure{peregrinations}, where it is $L_0$ that I obstinately kept favoring in the formulas --- cf. \S44, the definitive notation appears only in \S45 ---. The preference given to \deleted{$\unsure{\varepsilon}$} $-\rho\sigma^{-1}$ (denoted $\varepsilon_1$) over its inverse $-\sigma\rho^{-1}$ (denoted $\varepsilon_1^{-1}$) comes from the preference (justified by the facts) given to $\rho$ over $\sigma=\sigma_\infty$, in the loop formula written in the form
\[
[\alpha,\rho]=[\beta,\sigma]\varepsilon_1(\nu)
\]
instead of writing it
\[
[\beta,\sigma]=[\alpha,\rho]\varepsilon'_1(\nu)
\qquad(\varepsilon'_1=\varepsilon_1^{-1})
\]
So it is $\varepsilon_1$ that then \unsure{deserved} to have the simplest possible matrix form --- we took
\[
\varepsilon_1=\begin{pmatrix}1&0\\1&1\end{pmatrix}\ \ldots
\]

2) In terms of the systems
\[
\tag{64}
\diagram{\begin{array}{ccc}
\widetilde\ell'_\infty,&\widetilde\ell'_1,&\widetilde\ell'_0\\
\Vert&\Vert&\Vert\\
\widetilde\rho^{-1}&\widetilde\sigma&\widetilde\varepsilon_0
\end{array}}
\quad\text{satisfying}\quad
\widetilde\ell'_\infty\cdot\widetilde\ell'_1\cdot\widetilde\ell'_0=1
\]
(i.e. $\widetilde\rho^{-1}\cdot\widetilde\sigma\cdot\widetilde\varepsilon_{\unsure{0}}\unsure{=}1$ i.e. $\widetilde\varepsilon_0=\widetilde\sigma^{-1}\widetilde\rho$), the group $\unsure{\mathfrak S'_3|\mathcal G'_3}$ operates on them (cf. \S42 I) by cyclic permutations, and by the change into $(\widetilde\varepsilon_0^{-1},\widetilde\sigma^{-1},(\widetilde\rho^{-1})^{-1}=\widetilde\rho)$, \unsure{which} in its turn by cyclic permutation gives $(\widetilde\sigma^{-1},\widetilde\rho,\widetilde\varepsilon_0^{-1})$ ---

\mspage{365}

% Printed archive footer outside the manuscript frame:
% 365
% https://grothendieck.umontpellier.fr
% Université de Montpellier, UMR CNRS 5149, Case courrier 051, Place Eugène Bataillon
% 34095 Montpellier cedex 5 - France

\scanpage{367}

the passage from $(\widetilde\rho^{-1},\widetilde\sigma,\widetilde\varepsilon_0)$ to $(\widetilde\sigma^{-1},\widetilde\rho,\widetilde\varepsilon_0^{-1})$ then corresponds to the interchange $(\rho,\sigma)\mapsto(\sigma,\rho)$. How \deleted{\unsure{to apply some to|by applying to} $(\widetilde\rho^{-1},\widetilde\sigma,\widetilde\varepsilon_0)$} to interpret the transformation $(\rho,\sigma)\mapsto(\rho^{-1},\sigma^{-1})$, i.e. $(\widetilde\rho^{-1},\widetilde\sigma,\widetilde\varepsilon_0)\mapsto(\widetilde\rho,\widetilde\sigma^{-1},\widetilde\varepsilon_1)$ ?

Let us interpret the triples as homomorphisms.
\diagram{\begin{tikzcd}[ampersand replacement=\&,column sep=large,row sep=small]
\Pi_{0,3}\arrow[r,"\unsure{\varphi|\psi}"]\&\Gamma PG\\
\ell_0\arrow[r,mapsto]\&\ell'_0\\
\ell_1\arrow[r,mapsto]\&\ell'_1\\
\ell_\infty\arrow[r,mapsto]\&\ell'_\infty
\end{tikzcd}}
we see that the change $(\rho,\sigma)\mapsto(\sigma,\rho)$ corresponds to the composition of $\unsure{\varphi}$ with the automorphism of $\Pi_{0,3}$ such that
\diagram{\begin{tikzcd}[ampersand replacement=\&,column sep=small]
\ell_0\arrow[r,mapsto]\&\ell_1^{-1}\&\ell_1\arrow[r,mapsto]\&\ell_0^{-1}\&\text{thus }\ell_\infty\arrow[r,mapsto]\&\ell_\infty^{-1}
\end{tikzcd}}
i.e. with \uline{$\tau'_\infty$}. Likewise the transformation $(\rho,\sigma)\mapsto(\rho^{-1},\sigma^{-1})$ corresponds to composition with the automorphism
\diagram{\begin{tikzcd}[ampersand replacement=\&,column sep=small,row sep=small]
\ell_0\arrow[r,mapsto]\&\ell_0^{-1}\&\text{hence }\ell_\infty\arrow[r,mapsto]\&\ell_0\ell_1\\
\ell_1\arrow[r,mapsto]\&\ell_1^{-1}\&\&{}
\end{tikzcd}},
which is none other than \uline{$\tau_\infty$}$=\sigma_\infty\tau'_\infty$.

Thus the group of four transformations \deleted{\unsure{which operates|that we have}} on the set of pairs $(\rho,\sigma)$ that we considered in \uline{1$^\circ$)} is interpreted \deleted{\unsure{as}} \inserted{\deleted{by} as} the subgroup $\simeq\mathbb Z/2\mathbb Z\times\mathbb Z/2\mathbb Z$ generated by $\tau'_\infty,\tau_\infty$ in $\operatorname{Aut}_{\unsure{\mathrm{lac}|\mathrm{ext}}}(\Pi_{0,3})=\unsure{\mathfrak S'_{0,3}|\mathcal G_{0,3}}$, which \deleted{is} stabilizes $L_0\subset\Pi_{0,3}$ (thus its action on \unsure{the|this} third factor $\widetilde\varepsilon_0=\ell'_{\unsure{0}}$ will be to preserve it or replace it by its inverse, \deleted{\unsure{).}} which will have the virtue of \deleted{\unsure{it|it}} preserving its unipotent character, which we care about !).

\mspage{366}

% Printed archive footer outside the manuscript frame:
% 366
% https://grothendieck.umontpellier.fr
% Université de Montpellier, UMR CNRS 5149, Case courrier 051, Place Eugène Bataillon
% 34095 Montpellier cedex 5 - France

\scanpage{368}
\mspage{485}

\uline{$\overline{\mathrm{VIII}}$}\quad
\uline{Recapitulation on the equation}
\[
\tag{1}
[\alpha,\rho]=[\beta,\sigma]\varepsilon_1(\nu).
\]
The groups \deleted{$\unsure{\mathcal G_{\rho,\sigma}^{*0}}$}, $S\mathcal G_{\rho,\sigma}^{*}$ and its variants.

The essential result we have arrived at in (\uline{$\overline{\mathrm{III}}$}), is the possibility of parametrizing the solutions of such an equation \inserted{in $\alpha,\beta,\nu$} (under suitable hypotheses on $\rho,\sigma$) by a group $B'_0\times T_\rho\times T_\sigma$, where $B'_0$ is a ``\unsure{pre-Borel|pri-Borel}'' in $G$ (taking, for a suitable basis, the form
\[
B'_0=\left\{\begin{pmatrix}\mu&\lambda\\0&1\end{pmatrix}\;\middle|\;
\begin{array}{l}\mu\in\Gamma \mathbb G_m\\\lambda\in\Gamma \mathbb G_a\end{array}\right.
\]
and $T_\rho,T_\sigma$ are the \unsure{centralizers|stabilizers} of the regular sections $\rho,\sigma$ of $G$, by the morphism
\[
\tag{2}
\left\{\begin{gathered}
\diagram{\begin{tikzcd}[ampersand replacement=\&,column sep=large,row sep=small]
B'_0\times T_\rho\times T_\sigma\arrow[r]\&\mathcal G_{\rho,\sigma}\\
(\underline u,\underline a,\underline b)\arrow[r,mapsto]\&(\underline u\underline a^{-1},\underline u\underline b^{-1},\unsure{\nu|\delta}(\underline u))
\end{tikzcd}}
\end{gathered}\right.
\]
(where the fact that $\nu(\underline u)=\det(\underline u)-1$ is not of capital importance for the moment).

\margin{We now start from $\rho,\sigma,B'_0,\varepsilon_1$ (hence $L_1$) and we ask the question under what conditions, one can have a parametrization (2).}

\deleted{For} \inserted{To say} that this morphism \unsure{from the first|from the first ones} \unsure{into} $G\times G\times E'$ \deleted{gives} \inserted{maps} indeed into the solutions of (1), \deleted{it is necessary and sufficient} \inserted{\unsure{means}} also that we have identically
\[
\underline u\rho(\underline u)^{-1}=\underline u\sigma(u)^{-1}\varepsilon_1(\underline\nu(u))
\]
which can also be written, canceling $\unsure{\underline u|=}$ and taking inverses
\[
\tag{3}
\rho(\underline u)=\varepsilon_1(-\underline\nu(u))\sigma(\underline u).
\]
The possibility of finding $\underline\nu:B'_0\longrightarrow \mathbb G_a$ (morphism of schemes) \unsure{will give|would give}, via (2), a morphism $B'_0\times T_\rho\times T_\sigma\longrightarrow\mathcal G_{\rho,\sigma}\overset{\mathrm{dfn}}{=}$ scheme

\mspage{367}

% Printed archive footer outside the manuscript frame:
% 367
% https://grothendieck.umontpellier.fr
% Université de Montpellier, UMR CNRS 5149, Case courrier 051, Place Eugène Bataillon
% 34095 Montpellier cedex 5 - France

\scanpage{369}

of solutions of (1), is therefore equivalent to that
\deleted{\unsure{that}} \deleted{for every} $\sigma$ and $\rho$
\deleted{\unsure{have}} the same action
on $B'_0$, \uline{modulo} multiplication on the \unsure{left} by
\deleted{$\unsure{L_1}$} the subgroup \deleted{$\unsure{L_1}$}
\inserted{\unsure{unipoten\textsuperscript{t}}} with one parameter $L_1$ which
enters into the very writing of the equation (1) we are studying.
This is hard to conceive unless one admits that
$L_1\subset\rho(B'_0)\overset{\mathrm{df}}=B'_1$, which moreover implies
$\sigma(B'_0)\subset B'_1$ (and by a known argument $\sigma(B'_0)=B'_1$)
and $L_1=(B'_1)_u$. Conversely, if we have
$\rho(B'_0)=\sigma(B'_0)$, \deleted{\unsure{then for every}}
\inserted{i.e. $\sigma^{-1}$}\deleted{$\unsure{\rho(-)}$}$\rho$ normalizes $B'_0$,
\deleted{then} i.e. $\sigma^{-1}\rho$ a section of $B_0$, then
$\sigma(u)=\sigma u\sigma^{-1}$ and $\rho(u)=\rho u\rho^{-1}$ are two sections
of $B'_1$ which have the same determinant -- \unsure{hence abelianizatio\textsuperscript{n}},
hence which do coincide modulo $(B'_1)_u$.

Thus we see that the essential property for obtaining solutions of (1) via (2),
is to have
\[
(4)\qquad (\sigma^{-1}\rho)(B'_0)=B'_0\qquad\text{i.e.}\quad
\sigma^{-1}\rho\in\Gamma B_0\, ..
\]
\margin{\uline{N.B.} This condition \unsure{(if} $\sigma^{-1}\rho$ regular,
\unsure{i.e.} $\rho$ and $\sigma$ \unsure{are not homothetic}
\unsure{in any fiber)}, practically determines $B_0$
-- \deleted{\unsure{determines}} \unsure{over a field, there is an}
\uline{\unsure{ambiguity}} \unsure{possible} \unsure{between|of order} \uline{two}
(the \uline{\unsure{two}} Borels containing \unsure{the element in question}).}

For uniqueness of the representation of a triple $(\alpha,\beta,\nu)$ in the form
$(\underline u\underline a^{-1},\underline u\underline b^{-1},\underline\nu(\underline u))$,
it is necessary \deleted{to assume} \inserted{and it} suffices that we have
\[
(5)\qquad T_\sigma\cap T_\rho\cap B'_0=\{1\},
\]
\margin{It suffices for example that we have $T_\rho\cap T_\sigma=\unsure{\mathbb \mathbb G_m}$
and $\mathbb \mathbb G_m\cap B'_0=\{1\}$.}
\deleted{which is also equivalent to \unsure{our having}
\[
(6)\qquad T_\sigma\cap\unsure{\rho} B'_0\unsure{=}
\]}
the indeterminacy is

\mspage{368}

% Printed archive footer outside the manuscript frame:
% 368
% https://grothendieck.umontpellier.fr
% Université de Montpellier, UMR CNRS 5149, Case courrier 051, Place Eugène Bataillon
% 34095 Montpellier cedex 5 - France

\scanpage{370}
\mspage{486}

in the form
\diagram{\begin{tikzcd}[ampersand replacement=\&]
(\underline u,\underline a,\underline b) \arrow[r,mapsto] \&
(\underline u l,\underline a l,\underline b l)
\end{tikzcd}},
with $l$ a section of $T_\rho\cap T_\sigma\cap B'_0$.
There remains the question of the surjectivity of (2) -- can every solution of (1)
indeed be written in the form
$(u a^{-1},\underline u b^{-1},\nu(\underline u))$ ?
It amounts to the same (if $\rho,\sigma$ are regular)
to say that \deleted{\unsure{it}} we have surjectivity in
\[
(6)\qquad
\diagram{\begin{tikzcd}[ampersand replacement=\&]
B'_0 \arrow[r] \& \mathcal X_{\rho,\sigma} \\
\underline u \arrow[r,mapsto] \&
([\underline u,\rho],[\underline u,\sigma],\underline\nu(\underline u))
\end{tikzcd}}
\]
where $\mathcal X_{\rho,\sigma}$ is the scheme of simultaneous solutions
\deleted{in} in \deleted{$\unsure{\underline\nu,}$}
$\zeta,\eta,\nu$ of
\[
(7)\qquad
\det\zeta=\det\eta=1,\quad
\operatorname{Tr}\zeta\rho=\operatorname{Tr}\rho,\quad
\operatorname{Tr}\eta\sigma=\operatorname{Tr}\sigma,\quad
\zeta=\eta\varepsilon_1(\nu).
\]

It was essentially to obtain this surjectivity that we had been led
in section (II) to make the draconian hypotheses
(pages\ 435',436) which are expressed, at the cost of an étale,
by the conditions \deleted{\unsure{that|for}}
\[
(8)\qquad
\left\{\begin{array}{ll}
\mathrm{a)} & \sigma^{-1}\rho\text{ a section of }(B_0)_u=L_0
\quad(\text{and not only of }B_0)\\
\mathrm{b)} & B_0\text{ and }B_1=\rho(B_0)=
\deleted{\unsure{\rho}}\inserted{\sigma}(B_0)
\text{ in ``general position'',}\\
&\text{i.e. }L_0\text{ and }L_1=\rho(L_0)=\sigma(L_0)
\text{ independent;}\\
&\text{(in each fiber)}
\end{array}\right.
\]
but it is not clear whether these conditions
(suggested by earlier calculations, and indeed giving rise to
\unsure{some|some} very pretty calculations) are really indispensable for the result of parametrizati-

\mspage{369}

% Printed archive footer outside the manuscript frame:
% 369
% https://grothendieck.umontpellier.fr
% Université de Montpellier, UMR CNRS 5149, Case courrier 051, Place Eugène Bataillon
% 34095 Montpellier cedex 5 - France

\scanpage{371}

In our context it is \inserted{\unsure{quite}} natural to assume
$\varepsilon_0=\unsure{\sigma^{-1}\rho|-\sigma^{-1}\rho}$ unipotent
(or at least \unsure{quasi-unipotent}, i.e. \deleted{\unsure{a par}}
an iterate $\varepsilon_0^n$ unipotent) but condition b) of (8), more subtle,
appears less imperative. The sequel will show whether this restriction is judicious.

In the case where one considers an equation
\[
(9)\qquad
\underbrace{[\alpha,\rho]}_{\zeta}=
\varepsilon\underbrace{[\beta,\sigma]}_{\unsure{\eta}}
\varepsilon_1(\nu)
\qquad\deleted{\unsure{\varepsilon}}
\]
with a parameter $\varepsilon$ (which can in particular take the important value $-1$)
$\in\Gamma\mu_2$, \unsure{one|where} we have for the first time in these calculations
destroyed the symmetry between $\rho$ and $\sigma$, by introducing the condition
$\operatorname{Tr}\sigma=0$ as a condition of stability of the scheme
\inserted{$\mathcal X_\sigma$} of the $\eta$ satisfying
\[
\det\eta=1,\qquad\operatorname{Tr}\eta\sigma=\operatorname{Tr}\sigma
\]
under $\mu_2$. The essential result we were after, was the isomorphism
\[
(10)\qquad
\diagram{\begin{tikzcd}[ampersand replacement=\&,column sep=large]
B'_0\times T_\rho\times N'_\sigma \arrow[r,"\sim"] \& \mathcal G_{\rho,\sigma}^{*}\\
(\underline u,\underline a,\underline b) \arrow[r,mapsto] \&
(\alpha=\underline u\underline a^{-1},\beta=\underline u\underline b^{-1},
\mu=\det\underline u,\varepsilon=\varepsilon(\underline b))
\end{tikzcd}}
\]
where $N'_\sigma\subset N_\sigma$ is now the normalizer of $T_\sigma$
-- defined by the condition
\[
(11).\qquad \underline b\,\sigma(\underline b)^{-1}=
\varepsilon(\underline b)\deleted{\unsure{\sigma}}
\qquad\text{i.e.}\quad\underline b(\sigma)=\varepsilon(\underline b)\sigma
\]

\mspage{370}

% Printed archive footer outside the manuscript frame:
% 370
% https://grothendieck.umontpellier.fr
% Université de Montpellier, UMR CNRS 5149, Case courrier 051, Place Eugène Bataillon
% 34095 Montpellier cedex 5 - France

\scanpage{372}
\mspage{487}

That the morphism (10) does map into $\mathcal G_{\rho,\sigma}^{*}$,
is trivial by definition of $N'_\sigma$ and via the analogous fact for (2)
-- without a hypothesis of the kind $\operatorname{Tr}\sigma=0$.
But the possibility of obtaining (\unsure{local\textsuperscript{ly}})
any $\varepsilon$ via a $\underline b\in\Gamma N'_\sigma$
-- or simply $\varepsilon=-1$, if 2 invertible) requires precisely
$\operatorname{Tr}\sigma=0$.

Of course, we are going to consider $\mathcal G_{\rho,\sigma}^{*}$ as a group scheme,
by transport of \unsure{structure} from the product group structure of
$B'_0\times T_\rho\times N'_\sigma$.
We single out in it the subgroup
\[
(12)\qquad \mathcal G_{\rho,\sigma}^{\unsure{\flat|b}*}(0)
\quad\text{corresponding to :}\quad T'_0\times T_\rho\times N'_\sigma
\]
where
\[
(13)\qquad T_0=B'_0\cap B'_1
\]
is the distinguished maximal torus of $B'_0$, which
\deleted{for our} the canonical pinning is described as
\[
(14)\qquad T_0=\left\{\left(\begin{array}{cc}\mu&0\\0&1\end{array}\right.\unsure{\big\}|\big)}
\,\middle|\,\mu\in\Gamma\mathbb \mathbb G_m\right\}.
\]
We likewise set
\[
(15)\qquad \mathcal G_{\rho,\sigma}^{\unsure{\flat|b}}(0)
=\mathcal G_{\rho,\sigma}\cap\mathcal G_{\rho,\sigma}^{*\unsure{\flat|b}}(0),
\quad\text{corresponding to : }T'_0\times T_\rho\times T_\sigma
\]

\mspage{371}

% Printed archive footer outside the manuscript frame:
% 371
% https://grothendieck.umontpellier.fr
% Université de Montpellier, UMR CNRS 5149, Case courrier 051, Place Eugène Bataillon
% 34095 Montpellier cedex 5 - France

\scanpage{373}

We shall thus have semidirect product structures
\[
(16)\qquad
\mathcal G_{\rho,\sigma}\simeq
\mathcal G_{\rho,\sigma}^{\unsure{\flat|b}}(0)\cdot\Lambda_0,
\qquad
\mathcal G_{\rho,\sigma}^{*}\simeq
\mathcal G_{\rho,\sigma}^{*\unsure{\flat|b}}(0)\cdot\Lambda_0,
\]
where
\[
(17)\qquad
\left\{\begin{array}{l}
\displaystyle\Lambda_0\subset\mathcal G_{\rho,\sigma},\quad
\Lambda_0=\left\{\left(
\alpha=\underset{\substack{\Vert\\\varepsilon_0^{\unsure{n|u}}}}{\lambda},
\beta=\underset{\substack{\Vert\\\varepsilon_0^{\unsure{n|u}}}}{\lambda},
\mu=1\right)\,\middle|\,\lambda\in\Gamma L_0\right\}\\
\hfill(\unsure{n|u}\in\Gamma\mathbb \mathbb G_a)\\[4pt]
\text{corresponding to:}\quad L_0\times\{1\}\times\{1\}
\subset B'_0\times T_\rho\times T_\sigma.
\end{array}\right.
\]

We shall finally consider another subgroup \inserted{scheme} of
$\mathcal G_{\rho,\sigma}^{*}$, by the conditions
\[
(18)\qquad
\left\{\begin{array}{l}
\det\alpha\in\unsure{\Gamma\mu_2}\quad\text{i.e. }(\det\alpha)^2=1\\
\det\beta=1
\end{array}\right.
\qquad\left|\quad\text{whence }S\mathcal G_{\rho,\sigma}^{\unsure{*|\flat}}
\subset\mathcal G_{\rho,\sigma}^{*}\right.
\]

\deleted{Let then} \inserted{\unsure{Remark}}
\deleted{$\unsure{\varepsilon'\longleftarrow(\alpha,\beta),|\varepsilon'(\alpha,\beta),}$}
$S\mathcal G_{\rho,\sigma}^{*}$ can be interpreted as the scheme of systems
\[
(\alpha,\beta,\mu,\varepsilon,\varepsilon')
\]
satisfying
\[
(19)\qquad [\alpha,\rho]=\varepsilon[\beta,\sigma]\varepsilon_1(\mu-1),
\quad\det\alpha=\varepsilon',\quad
\inserted{\det\beta=1,}\quad
\underline{\varepsilon^2=\varepsilon'^2=1}
\]
\margin{N.B. $\varepsilon^2=1$ already follows from the first relation.}

In terms of the systems
\[
(\underline u,\underline a,\underline b)\in\Gamma B'_0\times T_\rho\times N'_\sigma
\]
the equations (19) (of which the \unsure{first|first} is automatic), reduce to:
\[
(20)\qquad\det\underline u=\det\underline b=\varepsilon'\det\underline a,
\qquad\varepsilon'^2=1,
\]

\mspage{372}

% Printed archive footer outside the manuscript frame:
% 372
% https://grothendieck.umontpellier.fr
% Université de Montpellier, UMR CNRS 5149, Case courrier 051, Place Eugène Bataillon
% 34095 Montpellier cedex 5 - France

\scanpage{374}
\mspage{488}
which \unsure{show} that $S\mathcal{G}_{\rho,\sigma}^{*}$ is \inserted{indeed} a \uline{subgroup}
of $\mathcal{G}_{\rho,\sigma}^{*}$. This subgroup contains $\Lambda_0$
(\unsure{corresponding}: $L_0\times\unsure{\{1\}\times\{1\}|\{(1,1)\}}\subset B'_0\times T_\rho\times N'_\sigma$) --- which
corresponds to the case where \deleted{\unsure{we have}}
\[
\det\underline{u}=\det\underline{b}=\det\underline{a}=1\qquad(\varepsilon'=1),
\tag{21}
\]
and we have a $\unsure{\frac12}$ direct product structure
\[
S\mathcal{G}_{\rho,\sigma}^{*}\simeq S\mathcal{G}_{\rho,\sigma}^{*\unsure{\flat|b}}(0)\mathbin{.}\Lambda_0
\tag{22}
\]
where
\[
S\mathcal{G}_{\rho,\sigma}^{*}(0)=S\mathcal{G}_{\rho,\sigma}^{*}\cap\mathcal{G}_{\rho,\sigma}^{*\unsure{\flat|b}}(0)
\tag{23}
\]
corresponds to the \unsure{subgroup|subgroup} \deleted{of} $\deleted{\unsure{B'_0}}\unsure{T_0|T'_0}\times T_\rho\times N'_\sigma$
\deleted{of $B'_0\times T_\rho\times N_\sigma$} defined by the same equations
(20). Of course, we shall also have
\[
S\mathcal{G}_{\rho,\sigma},\quad S\mathcal{G}_{\rho,\sigma}^{\unsure{\flat|b}}(0)\quad(\text{without }*)
\tag{24}
\]
\deleted{as} \inserted{intersections} of $S\mathcal{G}_{\rho,\sigma}^{*}$, $S\mathcal{G}_{\rho,\sigma}^{\unsure{\flat|b}*}(0)$ with
$S\mathcal{G}_{\rho\sigma}$ --- deduced from the expression
of these in terms of $(B'_0,\unsure{T_0|T'_0}),T_\rho,N'_\sigma$
by replacing $N'_\sigma$ with $T_\sigma$.

To make explicit the structure of the $S\mathcal{G}_{\rho,\sigma}^{*}$,
we are reduced via (22) \uline{to studying the}
group
\fbox{\diagram{\begin{tikzcd}[ampersand replacement=\&]
(25)\quad S\mathcal{G}_{\rho,\sigma}^{*\unsure{\flat|b}}(0) \arrow[r, hook] \& \unsure{T_0|T'_0}\times T_\rho\times N'_\sigma
\end{tikzcd}}}
\mspage{373}

% Printed archive footer outside the manuscript frame:
% 373
% https://grothendieck.umontpellier.fr
% Université de Montpellier, UMR CNRS 5149, Case courrier 051, Place Eugène Bataillon
% 34095 Montpellier cedex 5 - France

\scanpage{375}
described by equations (20). But we have
\diagram{\begin{tikzcd}[ampersand replacement=\&]
T_0 \arrow[r,"\sim"] \& \mathbb{G}_m \\
\underline{u} \arrow[r,mapsto] \& \det\underline{u}
\end{tikzcd}},
which implies that the projection
\diagram{\begin{tikzcd}[ampersand replacement=\&,column sep=large]
S\mathcal{G}_{\rho,\sigma}^{*\unsure{\flat|b}}(0) \arrow[r] \& T_\rho\times N'_\sigma \\
(\underline{u}\underline{a}^{-1},\underline{u}\underline{b}^{-1},\det\underline{u},\varepsilon=\varepsilon(\underline{b})) \arrow[r] \arrow[d,phantom,"\wr" description] \& (\underline{a},\underline{b}) \\
(\underline{u},\underline{a},\underline{b}) \&
\end{tikzcd}}
is a monomorphism, and induces an isomorphism of
$S\mathcal{G}_{\rho,\sigma}^{*}(0)$ onto the \unsure{subgroup|subgroup} of $T_\rho\times N'_\sigma$
\unsure{defined|defined} by the equation
\[
\det\underline{b}=\varepsilon'\det\underline{a},\qquad \varepsilon'{}^2=1
\tag{21}
\]
we shall have
\[
(\det\underline{b})^2=(\det\underline{a})^2.
\]

In terms of the homomorphisms (\unsure{determinants}, in short)
\diagram{\begin{tikzcd}[ampersand replacement=\&]
T_\rho \arrow[r,"\delta_\rho"] \& \mathbb{G}_m, \& N'_\sigma \arrow[r,"\unsure{\delta_\sigma}"] \& \mathbb{G}_m
\end{tikzcd}}
this amounts to taking the inverse image of the
diagonal \inserted{of $\mathbb{G}_m\times \mathbb{G}_m$} under $\delta_\rho^2,\delta_\sigma^2$.

We see at once (since $\delta_\rho,\delta_\sigma$
are epimorphisms) that the homomorphism
\diagram{\begin{tikzcd}[ampersand replacement=\&,column sep=large]
(22)\quad S\mathcal{G}_{\rho,\sigma}^{*\unsure{\flat|b}}(0) \arrow[r] \& \mu_2\times\mu_2 \\
(\alpha,\beta,\mu,\varepsilon,\varepsilon') \arrow[r,mapsto] \& (\varepsilon,\varepsilon') \\
\text{i.e.}\quad(\underline{u},\underline{a},\underline{b}) \arrow[r,mapsto] \& \left(\underline{\varepsilon}(\underline{b}),\underbrace{\det\underline{b}/\det\underline{a}}_{\det(\underline{b}\underline{a}^{-1})}\right)
\end{tikzcd}}
is an epimorphism, and \deleted{the kernel} thus gives
rise to an exact sequence
\mspage{374}

% Printed archive footer outside the manuscript frame:
% 374
% https://grothendieck.umontpellier.fr
% Université de Montpellier, UMR CNRS 5149, Case courrier 051, Place Eugène Bataillon
% 34095 Montpellier cedex 5 - France

\scanpage{376}
\mspage{489}
\diagram{\begin{tikzcd}[ampersand replacement=\&,column sep=small]
(23)\quad1 \arrow[r] \& S\mathcal{G}_{\rho,\sigma}^{*\circ\unsure{\flat|b}}(0) \arrow[r] \& S\mathcal{G}_{\rho,\sigma}^{*\unsure{\flat|b}}(0) \arrow[r] \& \mu_2\times\mu_2 \arrow[r] \& 1
\end{tikzcd}}
where $S\mathcal{G}_{\rho,\sigma}^{\unsure{\circ\flat|\circ}}(0)$ is isomorphic to the subscheme of
$T_\rho\times T_\sigma$ defined by the equation
\[
\det\underline{a}=\det\underline{b}
\tag{24}
\]
and thus in turn breaks down into
\diagram{\begin{tikzcd}[ampersand replacement=\&,column sep=small]
(25)\quad1 \arrow[r] \& ST_\rho^{\unsure{\flat|b}}\unsure{\times} ST_\sigma^{\unsure{\flat|b}} \arrow[r] \& S\mathcal{G}_{\rho,\sigma}^{\circ\unsure{\flat|b}}(0) \arrow[r] \& \mathbb{G}_m \arrow[r] \& 1 \\
\& \& \deleted{\unsure{(\underline{u},|\underline{u}}}(\underline{a},\underline{b}) \arrow[r,mapsto] \& \deleted{\det\underline{u}}\det\underline{a}=\det \&
\end{tikzcd}}
where we have set
\[
\left\{
\begin{aligned}
ST_\rho^{\unsure{\flat|b}}&=\operatorname{Ker}\left(T_\rho\xrightarrow{\delta_\rho}\mathbb{G}_m\right)\\
ST_\sigma^{\unsure{\flat|b}}&=\operatorname{Ker}\left(T_\sigma\xrightarrow{\delta_\sigma}\mathbb{G}_m\right)
\end{aligned}
\right.
\tag{26}
\]

If for example $\rho$ \inserted{(\unsure{non-s|non-g})} is $\frac12$ simple, hence $T_\rho$ a
maximal torus, then $ST_\rho^{\unsure{\flat|b}}$ is itself a torus \inserted{(of \unsure{dimension 1})} ---
\unsure{ditto} for $\sigma$. Thus, \deleted{\illegible{1}} in the ``modular case''
\[
\rho=\begin{pmatrix}0&1\\-1&1\end{pmatrix},\qquad
\sigma=\begin{pmatrix}0&-1\\1&0\end{pmatrix},
\]
$ST_\rho^{\unsure{\flat|b}}$ is a
torus \deleted{of} except in characteristic 3, and $ST_\sigma^{\unsure{\flat|b}}$ a torus except in
characteristic 2. \uline{Thus}, \uline{except in characteristic \inserted{2,3}}, $S\mathcal{G}_{\rho,\sigma}^{\circ\unsure{\flat|b}}(0)$ is a
\uline{torus of dimension 3}.

To finish determining the structure
of $S\mathcal{G}_{\rho,\sigma}^{*}$ in terms of $\mathcal{G}_{\rho,\sigma}^{*\unsure{\circ}}(0)$ and
of the decomposition (22) as a semidirect product, we must \deleted{specify}\inserted{make explicit} the operation of
\mspage{375}

% Printed archive footer outside the manuscript frame:
% 375
% https://grothendieck.umontpellier.fr
% Université de Montpellier, UMR CNRS 5149, Case courrier 051, Place Eugène Bataillon
% 34095 Montpellier cedex 5 - France

\scanpage{377}
$S\mathcal{G}_{\rho,\sigma}^{*\circ}(0)$ on $\Lambda_0$, by \unsure{interpreting}
$S\mathcal{G}_{\rho,\sigma}^{*\circ}(0)$ as a subgroup of $T_\rho\times N'_\sigma$,
and $\Lambda_0$ as \inserted{$L_0$} $\simeq \mathbb{G}_{\unsure{a|0}}$ via $\lambda\mapsto\varepsilon_0(\lambda)$.

Now the relations
\[
\deleted{\underline{u}\unsure{\varepsilon_0(\lambda)}\underline{u}^{-1}=}
\]
\[
\underline{u}(\varepsilon_0(\lambda))\underline{u}^{-1}
 =\varepsilon_0(\det(\underline{u})\lambda),\qquad
\det(\underline{u})=\det(\underline{b})=\delta_\sigma(\underline{b})
\]
show that the operation takes place via
the composite character
\diagram{\begin{tikzcd}[ampersand replacement=\&]
(27)\quad S\mathcal{G}_{\rho,\sigma}^{*\unsure{\circ|\flat}}(0) \arrow[r] \& N'_\sigma \arrow[r,"\delta_\sigma"] \& \mathbb{G}_m
\end{tikzcd}}, OK.

For what follows, we shall need
a canonical homomorphism
\diagram{\begin{tikzcd}[ampersand replacement=\&,column sep=large,row sep=large]
(28)\quad\mathcal{G}_{\rho,\sigma}^{*} \arrow[rr,"\unsure{\theta|\Theta}"] \arrow[dr,shift right=0.5ex,"{\deleted{\begin{array}{cccc}(\unsure{\alpha},&\beta,&\mu,&\unsure{\varepsilon})\\\unsure{\Vert|=}&\unsure{\Vert|=}&&\\\underline{u}\underline{a}^{-1}&\underline{u}\underline{b}^{-1}&&\end{array}}}"'] \& \& G \\
\& B'_0\times T_\rho\times\unsure{N'_\sigma|T'_\sigma} \arrow[ul,shift right=0.5ex] \arrow[ur,dashed] \&
\end{tikzcd}}
which identifies with the \uline{first} \uline{projection}
\diagram{\begin{tikzcd}[ampersand replacement=\&]
(\underline{u},\underline{a},\underline{b}) \arrow[r,mapsto] \& \underline{u}
\end{tikzcd}}
of the product, but \unsure{considered|considered} as
a morphism into $G$ thanks to the inclusion
\diagram{\begin{tikzcd}[ampersand replacement=\&]
B'_0 \arrow[r,phantom,"{\unsure{\hookrightarrow|\subset}}" description] \& G.
\end{tikzcd}}
We shall also \unsure{often} set
\unsure{(}
\[
\underline{u}=\unsure{\theta|\Theta}(\alpha,\beta,\mu,\varepsilon),
\]
\unsure{thus|whence} \deleted{if $\det$} we shall have
\mspage{376}

% Printed archive footer outside the manuscript frame:
% 376
% https://grothendieck.umontpellier.fr
% Université de Montpellier, UMR CNRS 5149, Case courrier 051, Place Eugène Bataillon
% 34095 Montpellier cedex 5 - France
